% Generated from ../chapters/23-objections-and-boundaries.md
\chapter{Chapter 23 --- Objections, Boundaries, and What Would Prove the Framework Wrong}\label{chapter-23-objections-boundaries-and-what-would-prove-the-framework-wrong}

A serious reader will raise objections. They should be addressed after the framework has been understood, not inserted defensively into every chapter.

\section{``These are ordinary benefits of exercise and stress reduction''}\label{these-are-ordinary-benefits-of-exercise-and-stress-reduction}

Perhaps. Improved posture, pain, sleep, and mood may be explained by known mechanisms without a broader theory. The stronger framework earns its place only if it predicts durable multiscale changes or improved future adaptability beyond conventional explanations.

\section{``Adult tissues lack the progenitors required for regeneration''}\label{adult-tissues-lack-the-progenitors-required-for-regeneration}

This is decisive in some cases but not universally. Mature-cell plasticity, transdifferentiation, and axonal regrowth show alternative pathways. The question must be tissue-specific.

\section{``Natural practice cannot reproduce targeted molecular therapy''}\label{natural-practice-cannot-reproduce-targeted-molecular-therapy}

Correct as an evidence statement: no practice has been shown to reproduce OSK reprogramming in human retinal ganglion cells. Incorrect as an impossibility proof: organism-level signals do alter chromatin, transcription, and cell state. Magnitude and specificity remain open.

\section{``The framework underestimates cancer''}\label{the-framework-underestimates-cancer}

Cancer is one of its strongest boundaries. Greater growth and plasticity are not automatically healthy. Any endogenous rejuvenation theory must explain how identity and tumour suppression are preserved.

\section{``The personal observations are confirmation bias''}\label{the-personal-observations-are-confirmation-bias}

They may be. This is why prospective measurement, blinded analysis, and replication are necessary.

\section{``No impossibility proof does not make the hypothesis true''}\label{no-impossibility-proof-does-not-make-the-hypothesis-true}

Agreed. Survival against objections keeps a hypothesis open; prediction and evidence make it credible.

The framework would be weakened if long-term practitioners show no greater structural or adaptive change than matched controls; if measured changes remain within normal variability; if improved state does not improve future adaptation; or if the apparent effects disappear under blinded measurement.

The eye remains a boundary case and should stay visible. A framework that explains only easy examples while hiding hard failures is not useful.
